\documentclass[11pt]{article}
\usepackage[margin=0.9in]{geometry}
\usepackage{graphicx}
\usepackage{caption}
\usepackage{hyperref}
\usepackage{float}
\usepackage{parskip}
\usepackage{booktabs}
\usepackage{array}

\hypersetup{colorlinks=false, hidelinks}
\setlength{\parskip}{3pt plus 1pt}
\captionsetup{font=small, skip=3pt, labelfont=bf}

\title{\textbf{GreenStep Employee Sustainability Challenge}\\[0.4em]
\large Source Code Documentation}
\author{Oli Gurmessa \and Hadi Shaar \and Naisha Srivastav \and Nathanael Agbemadon}
\date{CISC480 --- Senior Capstone\\[0.3em]
Instructor: Dr.\ Jimenez \quad\textbullet\quad Spring 2026\\[0.3em]
Submitted: May 22, 2026}

\begin{document}

\maketitle
\tableofcontents
\newpage

% ─────────────────────────────────────────────────────────────────────────────
\section{Overview and Tooling}
% ─────────────────────────────────────────────────────────────────────────────

This document describes the code-construction guidelines for the GreenStep
codebase: project structure, naming, imports, the server/client component
model, the API route pattern, commenting, formatting, and version-control
practice. The conventions follow the Google TypeScript Style Guide~\cite{gts}
adapted to a Next.js~15 + TypeScript project. Every example below cites a real
file in the repository so the rule is traceable to the code.

Two tools enforce quality, and both must pass before a pull request is merged:

\begin{itemize}
    \item \textbf{Prettier}~\cite{prettier} for formatting, configured by
    \texttt{.prettierrc} at the repository root.
    \item \textbf{ESLint}~\cite{eslint} for linting, configured by
    \texttt{eslint.config.mjs} (the Next.js flat-config preset
    \texttt{next/core-web-vitals} + \texttt{next/typescript}), run via
    \texttt{npm run lint}.
\end{itemize}

The continuous-integration gate is \texttt{npm run build} and
\texttt{npm run lint}. Automated tests are not yet configured for the project;
test plans and cases are maintained in the QA Activities document and listed
as future work in the Project Report.

% ─────────────────────────────────────────────────────────────────────────────
\section{Project / Folder Structure}
% ─────────────────────────────────────────────────────────────────────────────

The repository follows the Next.js App Router layout:

\begin{table}[H]
\centering
\caption{Top-level source directories and their responsibilities.}
\label{tab:struct}
\begin{tabular}{p{3.6cm} p{9.6cm}}
\toprule
Path & Responsibility \\
\midrule
\texttt{app/} & Routed pages and layouts (Presentation). Route groups
\texttt{(auth)/} and \texttt{(dashboard)/} organize related screens. \\
\texttt{app/api/} & Server-side API route handlers (\texttt{route.ts} files). \\
\texttt{app/actions/} & Server actions (e.g.\ \texttt{challenge.ts}). \\
\texttt{lib/} & Framework-agnostic helpers: \texttt{db.ts} (Prisma singleton),
\texttt{authz.ts}, \texttt{constants.ts}, \texttt{global-challenge.ts},
\texttt{team-requests.ts}, \texttt{actions-data.ts}. \\
\texttt{components/} & Shared, reusable UI components imported across
pages. \\
\texttt{prisma/} & \texttt{schema.prisma} --- the database schema. \\
\texttt{scripts/} & Developer seed scripts (\texttt{seed-*.js}). \\
\texttt{middleware.ts} & Clerk route protection (root). \\
\bottomrule
\end{tabular}
\end{table}

% ─────────────────────────────────────────────────────────────────────────────
\section{Naming Conventions}
% ─────────────────────────────────────────────────────────────────────────────

\begin{table}[H]
\centering
\caption{Naming conventions with a real example from the codebase.}
\label{tab:naming}
\begin{tabular}{p{3.0cm} p{4.4cm} p{5.6cm}}
\toprule
Case & Used for & Example (file) \\
\midrule
\texttt{PascalCase} & Components, types, interfaces &
\texttt{Home} in \texttt{app/page.tsx} \\
\texttt{camelCase} & Functions, variables, params &
\texttt{createTeamRequest}, \texttt{reviewTeamRequest} in
\texttt{lib/team-requests.ts} \\
\texttt{UPPER\_CASE} & Module constants &
\texttt{GLOBAL\_CHALLENGE\_ID} in \texttt{lib/constants.ts} \\
\texttt{kebab-case} & Route folders &
\texttt{app/(dashboard)/log-action/}, \texttt{app/api/team/request/} \\
\bottomrule
\end{tabular}
\end{table}

A component file is named after its main export (e.g.\
\texttt{DashboardClient.tsx} exports \texttt{DashboardClient}). Vague names
such as \texttt{utils.ts} are avoided in favor of scoped names
(\texttt{team-requests.ts}, \texttt{global-challenge.ts}).

% ─────────────────────────────────────────────────────────────────────────────
\section{Imports and Module Conventions}
% ─────────────────────────────────────────────────────────────────────────────

The project uses the \texttt{@/} path alias, configured in
\texttt{tsconfig.json} as \texttt{"@/*": ["./*"]}, so modules are imported by
absolute path rather than long relative chains.

Database access goes through a single Prisma client singleton in
\texttt{lib/db.ts}, imported everywhere as \texttt{db}:

\begin{verbatim}
// lib/db.ts
export const db =
  globalForPrisma.prisma ??
  new PrismaClient({ log: [...] });
if (process.env.NODE_ENV !== "production") globalForPrisma.prisma = db;

// usage in any route or lib module:
import { db } from "@/lib/db";
\end{verbatim}

The singleton prevents exhausting database connections during development
hot-reload.

% ─────────────────────────────────────────────────────────────────────────────
\section{Server vs.\ Client Components}
% ─────────────────────────────────────────────────────────────────────────────

Components are \textbf{server components by default}. Data fetching, auth, and
Prisma access happen on the server; for example,
\texttt{app/(dashboard)/dashboard/page.tsx} authenticates and syncs the user
before rendering.

A component that needs browser APIs, state, or event handlers opts in with the
\texttt{"use client"} directive on its first line --- for example,
\texttt{app/(dashboard)/dashboard/DashboardClient.tsx}:

\begin{verbatim}
"use client";
import { useState, useEffect } from "react";
// ... interactive dashboard
\end{verbatim}

Server-only mutations are written as \textbf{server actions} with the
\texttt{"use server"} directive, e.g.\ \texttt{app/actions/challenge.ts}
(\texttt{createChallenge}, \texttt{joinChallenge}).

% ─────────────────────────────────────────────────────────────────────────────
\section{API Route Pattern}
% ─────────────────────────────────────────────────────────────────────────────

Every API route handler follows the same shape: a \texttt{try/catch} wrapping
the logic, an authentication guard, input validation, and
\texttt{NextResponse.json(...)} with an explicit HTTP status code. The pattern
below is taken from \texttt{app/api/actions/route.ts}:

\begin{verbatim}
export async function POST(req: Request) {
  try {
    const { userId } = await auth();
    if (!userId)
      return NextResponse.json({ error: "Unauthorized" }, { status: 401 });

    const { category, actionType } = await req.json();
    if (!category || !actionType)
      return NextResponse.json(
        { error: "category and actionType are required" },
        { status: 400 });

    const user = await db.user.findUnique({ where: { clerkId: userId } });
    if (!user)
      return NextResponse.json({ error: "User not found" }, { status: 404 });

    // ... perform work ...
    return NextResponse.json({ ...action, streak: newStreak });
  } catch (error) {
    console.error(error);
    return NextResponse.json(
      { error: "Failed to create action" }, { status: 500 });
  }
}
\end{verbatim}

Status codes are used consistently across handlers: \texttt{401} unauthorized,
\texttt{400} bad input, \texttt{404} not found, \texttt{409} conflict (e.g.\
the duplicate-action guard), and \texttt{500} for unexpected errors. The same
shape appears in \texttt{app/api/leaderboard/route.ts} and
\texttt{app/api/me/route.ts}.

% ─────────────────────────────────────────────────────────────────────────────
\section{Commenting}
% ─────────────────────────────────────────────────────────────────────────────

Exported functions and components are documented with JSDoc. For functions
with parameters, a multi-line block documents each \texttt{@param}; the example
below is from \texttt{app/(dashboard)/dashboard/useCountUp.ts}:

\begin{verbatim}
/**
 * Animates a numeric value toward a target over a fixed duration.
 * @param value The target value to count toward.
 * @param duration Animation length in milliseconds (default 900).
 * @returns The current animated value to render.
 */
export function useCountUp(value: number, duration = 900) { ... }
\end{verbatim}

Self-explanatory helpers use a single-line JSDoc (see
\texttt{app/(dashboard)/dashboard/helpers.ts} and \texttt{useCountUp.ts}).
Implementation notes use \texttt{//} line comments.

\textbf{Aspirational rule --- file-level header JSDoc.} New modules begin with
a file-level JSDoc header (\texttt{File}, \texttt{Author}, \texttt{Created},
\texttt{Description}); this is present in newer files such as
\texttt{lib/constants.ts} and the extracted dashboard modules
(\texttt{types.ts}, \texttt{helpers.ts}, \texttt{useCountUp.ts}). It is a goal for new files rather than a property of
every existing file, and is not retrofitted across the whole codebase.

% ─────────────────────────────────────────────────────────────────────────────
\section{Formatting (Prettier)}
% ─────────────────────────────────────────────────────────────────────────────

Formatting is handled by Prettier, configured in \texttt{.prettierrc} at the
repository root (Table~\ref{tab:prettier}).

\begin{table}[H]
\centering
\caption{Prettier settings (\texttt{.prettierrc}).}
\label{tab:prettier}
\begin{tabular}{p{4.5cm} p{4.5cm}}
\toprule
Setting & Value \\
\midrule
\texttt{printWidth} & 80 \\
\texttt{tabWidth} & 2 \\
\texttt{semi} & true \\
\texttt{singleQuote} / \texttt{jsxSingleQuote} & false (double quotes) \\
\texttt{trailingComma} & es5 \\
\texttt{bracketSpacing} & true \\
\texttt{arrowParens} & avoid \\
\bottomrule
\end{tabular}
\end{table}

% ─────────────────────────────────────────────────────────────────────────────
\section{Version Control}
% ─────────────────────────────────────────────────────────────────────────────

\begin{itemize}
    \item \textbf{Branches} are named \texttt{type/short-description}, e.g.\
    \texttt{feature/leaderboard}, \texttt{fix/null-check-userid},
    \texttt{chore/update-deps}.
    \item \textbf{Commit messages} use the imperative mood and describe one
    logical change (e.g.\ ``Add sign-in button''), following Conventional
    Commits~\cite{cc}.
    \item \textbf{Quality gate.} \texttt{npm run build} and
    \texttt{npm run lint} must pass before a pull request is ready. (No
    \texttt{test} script is defined yet; automated testing is tracked in the
    QA Activities document.)
    \item \textbf{No debugging artifacts} (stray \texttt{console.log} or
    commented-out code) are committed.
\end{itemize}

% ─────────────────────────────────────────────────────────────────────────────
\begin{thebibliography}{9}
% ─────────────────────────────────────────────────────────────────────────────

\bibitem{gts}
Google, ``Google TypeScript Style Guide,'' 2024. [Online]. Available:
\url{https://google.github.io/styleguide/tsguide.html}

\bibitem{prettier}
Prettier, ``Prettier Documentation,'' 2024. [Online]. Available:
\url{https://prettier.io/docs/en/}

\bibitem{eslint}
ESLint, ``ESLint Documentation,'' 2024. [Online]. Available:
\url{https://eslint.org/docs/latest/}

\bibitem{nextjs}
Vercel, ``Next.js Documentation,'' 2024. [Online]. Available:
\url{https://nextjs.org/docs}

\bibitem{cc}
Conventional Commits, ``Conventional Commits v1.0.0,'' 2023. [Online].
Available: \url{https://www.conventionalcommits.org/}

\end{thebibliography}

\end{document}
